% !TEX program = xelatex
% Standalone viewer for the AR(2) simulation paired-test table.
% Table float generated by make_ar2_sim_table.R -> ar2_sim_table_body.tex.
\documentclass[12pt]{article}

\usepackage{fontspec}
\usepackage{xeCJK}
\setCJKmainfont[AutoFakeSlant=.2, AutoFakeBold=3]{AR PL KaitiM Big5}

\usepackage[margin=1in]{geometry}
\usepackage{amsmath, amssymb}
\usepackage{booktabs}
\usepackage{array}
\usepackage{hyperref}
\hypersetup{colorlinks=true, linkcolor=blue!45!black, urlcolor=blue!45!black}

\title{\textbf{模擬章的不確定性：AR(2) 空間 hold-out 配對檢定}}
\author{方法 2（複製層配對檢定，$n=10$ 組資料）}
\date{2026 年 7 月}

\begin{document}
\maketitle

\section*{方法：以 10 組資料為獨立單位}

臭氧實證用的是\textbf{站層 bootstrap}，因為只有 $2$--$3$ 個 split、站是唯一近似
獨立的重抽單位。模擬則不同：每組模擬資料是一次\textbf{獨立的實驗複製}，所以
獨立單位是「資料組」。論文 Figure~1（settings 9--12）目前只畫 Brier-對-分位的
曲線、沒有不確定帶；本表為每個配對對照，取\textbf{逐組 Brier 差}
$d_r = \mathrm{BS}^A_r - \mathrm{BS}^B_r$（$r=1,\dots,10$），做配對 $t$（$95\%$ CI）
與 Wilcoxon signed-rank 檢定。這正是 \verb|ar2_paired_tests.R| 直接讀
\verb|scores.R| 已快取的 per-dataset \verb|brier.score| 陣列所算。

\bigskip
% Auto-generated by make_ar2_sim_table.R -- do not edit by hand.
\begin{table}[htbp]
\centering
\caption{AR(2) simulation, spatial hold-out arm (settings 9--12): paired
Brier differences over the $10$ data sets, $\Delta=\mathrm{BS}_A-\mathrm{BS}_B$
in units of $10^{-3}$ (negative favours the first-named model), with $95\%$
paired-$t$ confidence intervals. Two core contrasts: the temporal model
against its non-temporal counterpart, and AR(2) against AR(1). Every
interval covers $0$ except one starred cell (setting~11, $q{=}0.90$), which
favours the \emph{baseline}, is borderline ($p_W{=}0.05$), and does not
survive multiplicity across the $36$ comparisons. This is the quantitative
form of ``the temporal curves sit on top of their non-temporal counterparts.''}
\label{tab:ar2-sim-paired}
\small
\begin{tabular}{@{}ll cc@{}}
\toprule
Setting & $L$ quantile & AR(2) $-$ no-TS & AR(2) $-$ AR(1) \\
\midrule
9 (K{=}1 strong) & $0.90$ & $+0.00$ {\scriptsize $[-0.03,\,+0.04]$} & $-0.03$ {\scriptsize $[-0.06,\,+0.00]$} \\
 & $0.95$ & $-0.01$ {\scriptsize $[-0.03,\,+0.01]$} & $-0.02$ {\scriptsize $[-0.06,\,+0.02]$} \\
 & $0.98$ & $-0.01$ {\scriptsize $[-0.04,\,+0.01]$} & $-0.01$ {\scriptsize $[-0.03,\,+0.01]$} \\
\addlinespace
10 (K{=}1 weak) & $0.90$ & $+0.00$ {\scriptsize $[-0.02,\,+0.03]$} & $-0.01$ {\scriptsize $[-0.04,\,+0.02]$} \\
 & $0.95$ & $-0.00$ {\scriptsize $[-0.03,\,+0.02]$} & $-0.01$ {\scriptsize $[-0.03,\,+0.01]$} \\
 & $0.98$ & $-0.01$ {\scriptsize $[-0.03,\,+0.01]$} & $+0.00$ {\scriptsize $[-0.01,\,+0.01]$} \\
\addlinespace
11 (K{=}5 strong) & $0.90$ & $+0.93^{*}$ {\scriptsize $[+0.12,\,+1.75]$} & $+0.36$ {\scriptsize $[-0.51,\,+1.23]$} \\
 & $0.95$ & $+0.32$ {\scriptsize $[-0.33,\,+0.97]$} & $+0.41$ {\scriptsize $[-0.12,\,+0.94]$} \\
 & $0.98$ & $-0.00$ {\scriptsize $[-0.33,\,+0.32]$} & $-0.14$ {\scriptsize $[-0.33,\,+0.05]$} \\
\addlinespace
12 (K{=}5 weak) & $0.90$ & $-0.21$ {\scriptsize $[-1.16,\,+0.74]$} & $-0.48$ {\scriptsize $[-1.16,\,+0.20]$} \\
 & $0.95$ & $+0.13$ {\scriptsize $[-0.27,\,+0.52]$} & $+0.08$ {\scriptsize $[-0.34,\,+0.50]$} \\
 & $0.98$ & $+0.04$ {\scriptsize $[-0.26,\,+0.35]$} & $+0.13$ {\scriptsize $[-0.37,\,+0.63]$} \\
\bottomrule
\end{tabular}
\end{table}


\bigskip
\noindent\textbf{結論。}$36$ 個對照中，$35$ 個的 $95\%$ CI 都涵蓋 $0$；唯一星號
（setting~11, $q{=}0.90$）\textbf{偏向基準}（AR(2) 較\emph{差}），$p_W=0.05$ 為
邊界值，且在 $36$ 個比較的多重性下（$\alpha=0.05$ 期望約 $1.8$ 個假陽性）不足
採信。也就是說：在\textbf{空間內插}任務上，AR(2) 相對「無時序」與相對 AR(1) 都
\emph{測不出} Brier 差異——這與論文的敘述一致，但現在有配對檢定作為量化證據。

\medskip
\noindent\textbf{與 time-block arm 的對照。}值得強調：同一個 AR(2) 模型，在
\emph{時間外推}（time-block forecast）arm 上\emph{確實}出現方向正確、量級大一個
數量級的增益（見 \verb|block1_positive_control/analyze_threeway.R| 及論文
persistence 小節）。兩者並置，正好支撐論文的核心論點：\textbf{AR(2) 是預測
（外推）工具，不是內插工具}。空間 hold-out 測不到它，不是因為沒訊號，而是因為
內插任務裡當日橫斷面資訊已主導一切。

\medskip
\noindent\textbf{嵌入論文。}把 \verb|ar2_sim_table_body.tex| 複製到論文目錄，於
\S\,``Results on AR(2) comparison'' 以 \verb|% Auto-generated by make_ar2_sim_table.R -- do not edit by hand.
\begin{table}[htbp]
\centering
\caption{AR(2) simulation, spatial hold-out arm (settings 9--12): paired
Brier differences over the $10$ data sets, $\Delta=\mathrm{BS}_A-\mathrm{BS}_B$
in units of $10^{-3}$ (negative favours the first-named model), with $95\%$
paired-$t$ confidence intervals. Two core contrasts: the temporal model
against its non-temporal counterpart, and AR(2) against AR(1). Every
interval covers $0$ except one starred cell (setting~11, $q{=}0.90$), which
favours the \emph{baseline}, is borderline ($p_W{=}0.05$), and does not
survive multiplicity across the $36$ comparisons. This is the quantitative
form of ``the temporal curves sit on top of their non-temporal counterparts.''}
\label{tab:ar2-sim-paired}
\small
\begin{tabular}{@{}ll cc@{}}
\toprule
Setting & $L$ quantile & AR(2) $-$ no-TS & AR(2) $-$ AR(1) \\
\midrule
9 (K{=}1 strong) & $0.90$ & $+0.00$ {\scriptsize $[-0.03,\,+0.04]$} & $-0.03$ {\scriptsize $[-0.06,\,+0.00]$} \\
 & $0.95$ & $-0.01$ {\scriptsize $[-0.03,\,+0.01]$} & $-0.02$ {\scriptsize $[-0.06,\,+0.02]$} \\
 & $0.98$ & $-0.01$ {\scriptsize $[-0.04,\,+0.01]$} & $-0.01$ {\scriptsize $[-0.03,\,+0.01]$} \\
\addlinespace
10 (K{=}1 weak) & $0.90$ & $+0.00$ {\scriptsize $[-0.02,\,+0.03]$} & $-0.01$ {\scriptsize $[-0.04,\,+0.02]$} \\
 & $0.95$ & $-0.00$ {\scriptsize $[-0.03,\,+0.02]$} & $-0.01$ {\scriptsize $[-0.03,\,+0.01]$} \\
 & $0.98$ & $-0.01$ {\scriptsize $[-0.03,\,+0.01]$} & $+0.00$ {\scriptsize $[-0.01,\,+0.01]$} \\
\addlinespace
11 (K{=}5 strong) & $0.90$ & $+0.93^{*}$ {\scriptsize $[+0.12,\,+1.75]$} & $+0.36$ {\scriptsize $[-0.51,\,+1.23]$} \\
 & $0.95$ & $+0.32$ {\scriptsize $[-0.33,\,+0.97]$} & $+0.41$ {\scriptsize $[-0.12,\,+0.94]$} \\
 & $0.98$ & $-0.00$ {\scriptsize $[-0.33,\,+0.32]$} & $-0.14$ {\scriptsize $[-0.33,\,+0.05]$} \\
\addlinespace
12 (K{=}5 weak) & $0.90$ & $-0.21$ {\scriptsize $[-1.16,\,+0.74]$} & $-0.48$ {\scriptsize $[-1.16,\,+0.20]$} \\
 & $0.95$ & $+0.13$ {\scriptsize $[-0.27,\,+0.52]$} & $+0.08$ {\scriptsize $[-0.34,\,+0.50]$} \\
 & $0.98$ & $+0.04$ {\scriptsize $[-0.26,\,+0.35]$} & $+0.13$ {\scriptsize $[-0.37,\,+0.63]$} \\
\bottomrule
\end{tabular}
\end{table}
| 置於
Figure~1 之後，並在文中把「sit essentially on top」一句改寫為引用
Table~\ref{tab:ar2-sim-paired} 的區間結論。表格由資料重生，數字免手動維護。

\end{document}
